\begin{table}[t!]
\centering
\footnotesize
\caption{\textbf{Pitch estimation methods used in the computational analysis studies for Carnatic and Hindustani music.} The FTANet-Carnatic model is an output of this thesis, and it is presented in detail in Chapter~\ref{ch04_melody}.}
\label{tab:pitch-estimation}
\begin{tabular}{p{7.4cm} c c}
\toprule
\textbf{Method} & \textbf{for polyph. input} & \textbf{Typology} \\
\midrule
Harmonic Product Spectrum (HPS)~\citep{noll-hps} & \xmark & \textcolor{olive}{Heuristics} \\
Autocorrelation~\citep{rabiner-autocorrelation} & \xmark & \textcolor{olive}{Heuristics} \\
YIN~\citep{decheveigne-yin} & \xmark & \textcolor{olive}{Heuristics} \\
\citep{klapuri-multiple_f0} & (Multiple-$f_0$ est.) & \textcolor{olive}{Heuristics} \\
SWIPE~\citep{camacho-sawtooth} & \xmark & \textcolor{olive}{Heuristics} \\
\citep{rao-pitch} & \cmark & \textcolor{olive}{Heuristics} \\
Melodia~\citep{salamon-melodia} & \cmark & \textcolor{olive}{Heuristics} \\
CREPE~\citep{kim-crepe} & \xmark & \textcolor{cyan}{DL} \\
FTANet-Carnatic~\citep{plaja-pitch}* & \cmark & \textcolor{cyan}{DL} \\
\bottomrule
\end{tabular}
\end{table}
